
\chapterspace{-4}
\chapter{Geometría algebraica y~geometría~analítica}
\label{XII}
\marginpar{311}

\begin{center}
{Sra. M.\ \textsc{Raynaud}\footnote{A partir de notas inéditas de
A\ptbl\ Grothendieck.}}
\end{center}
\vspace*{1cm}

Siguiendo el procedimiento de~\cite{XII.10}, a todo esquema~$X$ localmente
de tipo finito sobre el cuerpo de los números complejos~$\CC$ se le asocia
un espacio analítico~$X^\an$
\label{indnot:lb}\oldindexnot{$X^\an$|hyperpage}%
cuyo conjunto subyacente es~$X(\CC)$.

En los números~\Ref{XII.2} y~\Ref{XII.3} de esta exposición damos un
\og diccionario\fg entre las propiedades usuales de~$X$ y de~$X^\an$, y
entre las propiedades de un morfismo $f\colon X\to Y$ y las del morfismo
asociado $f^\an\colon X^\an\to Y^\an$.
\label{indnot:lc}\oldindexnot{$f^\an$|hyperpage}%

Demostramos después que los teoremas de comparación entre haces coherentes
sobre~$X$ y sobre~$X^\an$, establecidos en~\cite[\No 12]{XII.10} para una
variedad proyectiva, siguen siendo válidos cuando~$X$ es un esquema propio.

Por último, en el número~\Ref{XII.5} demostramos la equivalencia entre la
categoría de los recubrimientos étales finitos de~$X$ y la categoría de los
recubrimientos étales finitos de~$X^\an$. Como añadido para el lector, damos
una nueva demostración del teorema de Grauert--Remmert~\cite{XII.6}, que
utiliza la resolución de singularidades~\cite{XII.8}.

\section{Espacio analítico asociado a un esquema}
\label{XII.1}
\marginpar{312}

Sea~$X$ un esquema localmente de tipo finito sobre~$\CC$. Sea~$\Phi$ el
funtor de la categoría de los espacios analíticos~\cite[\No 9]{XII.4} en la
categoría de los conjuntos que asocia a un espacio analítico~$\othercal{X}$
el conjunto de los morfismos de espacios anillados en $\CC$-álgebras
$\Hom_\CC(\othercal{X},X)$. Se tiene el teorema siguiente:

\begin{theoremedefinition}
\label{XII.1.1}
El funtor~$\Phi$ es representable por un espacio analítico~$X^\an$ y un
morfismo $\varphi\colon X^\an\to X$. Se dice que~$X^\an$ es el espacio
analítico asociado a~$X$.
\index{espacio analítico asociado|hyperpage}%

Si~$|X^\an|$ es el conjunto subyacente a~$X^\an$, entonces~$\varphi$ induce
una biyección de~$|X^\an|$ sobre el conjunto~$X(\CC)$ de los puntos de~$X$
con valores en~$\CC$. Además, para cada punto~$x$ de~$X^\an$, el morfismo
$$
\varphi_x\colon\cal{O}_{X,\varphi(x)}\to\cal{O}_{X^\an,x},
$$
que es necesariamente local, induce al pasar a las completaciones un
isomorfismo
$$
\widehat{\varphi}_x\colon\widehat{\cal{O}}_{X,\varphi(x)}\isomto
\widehat{\cal{O}}_{X^\an,x}
\ifthenelse{\boolean{orig}}
{,}
{.}%
$$
En particular, el morfismo~$\varphi$ es plano.
\end{theoremedefinition}

Obsérvese que el hecho de que~$\varphi$ induzca una biyección de~$X^\an$
sobre~$X(\CC)$ resulta de la propiedad universal de~$X^\an$. Por otra parte,
se tienen las afirmaciones siguientes:

\begin{enumerate}
\item[a)] Si el teorema es cierto para un esquema~$Y$, también lo es para
todo subesquema~$X$ de~$Y$. Supongamos primero que~$X$ es un subesquema
abierto de~$Y$; si $\psi\colon Y^\an\to Y$ es el morfismo canónico,
$\psi^{-1}(X)$ es un abierto
\marginpar{313}
de~$Y^\an$
\ifthenelse{\boolean{orig}}
{al que se dota}
{al que se dota}%
~de la estructura de espacio analítico inducida por la de~$Y^\an$. Como todo
morfismo de un espacio analítico~$\othercal{X}$ en~$X$ se factoriza a través
de~$Y^\an$ por la propiedad universal de este último, y por tanto a través
de~$X^\an$, que es el producto fibrado~$Y^\an\times_YX$, resulta que~$X^\an$
es el espacio analítico asociado a~$X$. Por último, la afirmación relativa a
los~$\varphi_x$ es evidente.

Basta ahora considerar el caso en que~$X$ es un subesquema cerrado de~$Y$.
Sea~$I$ el $\cal{O}_Y$-ideal coherente que define~$X$; entonces
$I\cdot\cal{O}_{Y^\an}$ es un haz coherente de ideales de~$\cal{O}_{Y^\an}$
que define un subespacio analítico cerrado~$X^\an$ de~$Y^\an$; como en el
caso de un subesquema abierto, se ve que~$X^\an$ es el espacio analítico
asociado a~$X$. Sea $\varphi\colon X^\an\to X$ el morfismo canónico. Para
todo punto~$x$ de~$X^\an$, el morfismo~$\varphi_x$ no es sino el morfismo
$$
\cal{O}_{Y,\psi(x)}/I_{\psi(x)}\to
\cal{O}_{Y^\an,x}/I_{\psi(x)}\cdot\cal{O}_{Y^\an,x}
$$
inducido por~$\psi_x$; su completación
$$
\widehat{\varphi}_x\colon
\widehat{\cal{O}}_{Y,\psi(x)}/I_{\psi(x)}\cdot\widehat{\cal{O}}_{Y,\psi(x)}
\to\widehat{\cal{O}}_{Y^\an,x}/I_{\psi(x)}\cdot\widehat{\cal{O}}_{Y^\an,x}
$$
es un isomorfismo porque~$\widehat\psi_x$ lo es, lo que demuestra~a).

\item[b)] Si se tienen dos $\CC$-esquemas~$X_1$ y~$X_2$ tales que existen
$X_1^\an$ y~$X_2^\an$, entonces también existe~$(X_1\times X_2)^\an$. Sean,
en efecto, $\varphi_1\colon X_1^\an\to X_1$ y
$\varphi_2\colon X_2^\an\to X_2$ los morfismos canónicos, y sean~$p_1$
y~$p_2$ las dos proyecciones de~$X_1^\an\times X_2^\an$. Se deduce
formalmente de EGA~I~1.8.1 que~$X_1\times X_2$ es el producto de~$X_1$ y~$X_2$
en la categoría de los espacios localmente anillados; de ello resulta que
los morfismos~$\varphi_1\cdot p_1$ y~$\varphi_2\cdot p_2$ definen un
\marginpar{314}
morfismo $\varphi\colon X_1^\an\times X_2^\an\to X_1\times X_2$, y que el
par~$(X_1^\an\times X_2^\an,\varphi)$ representa el funtor
$\othercal{X}\mto\Hom_\CC(\othercal{X},X_1\times X_2)$.

\item[c)] Si~$\othercal{E}^1$ denota el espacio afín de dimensión~1, es
decir, el espacio topológico~$\CC$ provisto del haz de funciones holomorfas,
el funtor%
\ifthenelse{\boolean{orig}}
{}
{{\renewcommand{\thefootnote}{*}\addtocounter{footnote}{-1}\footnote{\lcrochetbf Añadido en 2003: $E^1_\CC$ denota la recta afín (algebraica) sobre~$\CC$.\rcrochetbf}}}
$\othercal{X}\mto\Hom_\CC(\othercal{X},E^1_\CC)$ es representable por
$\othercal{E}^1$, siendo el morfismo canónico
$\varphi\colon\othercal{E}^1\to E^1_\CC$ el morfismo evidente. En efecto,
darse un morfismo de un espacio analítico~$\othercal{X}$ en~$E^1_\CC$
equivale a darse un elemento de~$\Gamma(\othercal{X},\cal{O}_\othercal{X})$,
lo que equivale también a darse un morfismo de~$\othercal{X}$
en~$\othercal{E}^1$. Se tiene evidentemente una biyección
$|\othercal{E}^1|\simeq E^1(\CC)$ y, para cada punto
$x\in\othercal{E}^1$, el morfismo~$\widehat\varphi_x$ no es sino el
morfismo identidad de un anillo de series formales en una variable
sobre~$\CC$.
\end{enumerate}

\ifthenelse{\boolean{orig}}{}
{\enlargethispage{.5cm}}%
De b) y c) se deduce que el teorema es cierto para el espacio afín
$E^n_\CC$, $n\geq0$. Utilizando~a), se ve que también lo es para todo esquema
afín~$X$ localmente de tipo finito sobre~$\CC$. Si ya no se supone que~$X$
sea afín y si~$(X_i)$ es un recubrimiento de~$X$ por abiertos afines, se
deduce de la propiedad universal y de~a) que los~$X_i^\an$ se pegan y definen
así el espacio analítico~$X^\an$ asociado a~$X$.

\setcounter{subsection}{1}% Preserve the source-visible paragraph number XII.1.2.
\subsection{}
\label{XII.1.2}
Sea $f\colon X\to Y$ un morfismo de $\CC$-esquemas localmente de tipo finito.
Si $\varphi\colon X^\an\to X$ y $\psi\colon Y^\an\to Y$ son los morfismos
canónicos, de la propiedad universal de~$Y^\an$ resulta que existe un único
morfismo $f^\an\colon X^\an\to Y^\an$ tal que el diagrama
$$
\xymatrix{
X^\an \ar[r]\ar[d]_-{f^\an}& X\ar[d]^-f\\ Y^\an \ar[r]& Y
}
$$
sea
\marginpar{315}
conmutativo. Se ha definido, pues, un funtor~$\Phi$ de la categoría de los
$\CC$-esquemas localmente de tipo finito en la categoría de los espacios
analíticos.

El funtor~$\Phi$ conmuta con los límites proyectivos finitos. Basta, en
efecto, ver que~$\Phi$ conmuta con los productos fibrados. Ahora bien, si
$X$, $Y$ y~$Z$ son esquemas localmente de tipo finito sobre~$\CC$, del hecho
de que~$X\times_ZY$ sea el producto fibrado de~$X$ y~$Y$ sobre~$Z$ en la
categoría de los espacios localmente anillados resulta que
$X^\an\times_{Z^\an}Y^\an$ satisface la propiedad universal que caracteriza
a~$(X\times_ZY)^\an$.

\subsection{}
\label{XII.1.3}
Sean~$X$ un $\CC$-esquema localmente de tipo finito, $X^\an$ el espacio
analítico asociado y $\varphi\colon X^\an\to X$ el morfismo canónico. Si~$F$
es un $\cal{O}_X$-módulo, la imagen inversa
$\varphi^*F=F^\an$ es un haz de módulos sobre~$\cal{O}_{X^\an}$. Se define
así un funtor de la categoría de los $\cal{O}_X$-módulos en la categoría de
los módulos sobre~$X^\an$. Este funtor conmuta con los límites inductivos
(EGA~0~4.3.2). Como el haz~$\cal{O}_{X^\an}$ es coherente
\cite[\no18~\S 2~th\ptbl 2]{XII.4}, transforma haces coherentes en haces
coherentes (EGA~0~5.3.11). Se tiene además:

\begingroup
\setcounter{proposition}{0}% The source numbers this statement inside paragraph XII.1.3.
\renewcommand{\theproposition}{\thesubsection.\arabic{proposition}}
\renewcommand{\theHproposition}{\theHsubsection.\arabic{proposition}}
\begin{subproposition}
\label{XII.1.3.1}
El funtor que a un
\ifthenelse{\boolean{orig}}
{$\cal{O}_{X^\an}$-módulo}
{$\cal{O}_X$-módulo}%
~$F$ asocia su imagen inversa~$F^\an$ sobre~$X^\an$ es exacto, fiel y
conservativo.
\end{subproposition}
\endgroup

La exactitud resulta de que el morfismo $\varphi\colon X^\an\to X$ es plano
\ifthenelse{\boolean{orig}}
{\eqref{XIII.1.1}}
{\eqref{XII.1.1}}.%
\ Demostremos que el funtor~$F\mto F^\an$ es fiel. Habida cuenta de la
exactitud, basta demostrar que, si~$F^\an$ es nulo, también lo es~$F$. Ahora
bien, para todo punto~$x$ de~$X^\an$, se tiene entonces
$F_{\varphi(x)}\otimes_{\cal{O}_{X,\varphi(x)}}\cal{O}_{X^\an,x}=0$. Como el
morfismo $\cal{O}_{X,\varphi(x)}\to\cal{O}_{X^\an,x}$ es fielmente plano, se
tiene $F_{\varphi(x)}=0$ para todo punto cerrado~$\varphi(x)$ de~$X$ y, como
$X$ es de Jacobson (EGA~IV~10.4.8), esto implica que~$F$ es nulo.

El
\marginpar{316}
hecho de que el funtor~$F\mto F^\an$ sea conservativo resulta formalmente de
la exactitud y la fidelidad.
